% ============================================================
\subsection{Heaps and Priority Queues}
% ============================================================

A binary heap maintains the minimum of a collection under insertion and
deletion at logarithmic cost. Python's \texttt{heapq} operates on an ordinary
list interpreted as a heap, so no dedicated type is involved.

\begin{center}
\begin{tabular}{lll}
\toprule
\textbf{Operation} & \textbf{Call} & \textbf{Cost} \\
\midrule
Convert a list in place & \texttt{heapify(xs)}     & $O(n)$ \\
Insert                  & \texttt{heappush(h, x)}  & $O(\log n)$ \\
Remove the minimum      & \texttt{heappop(h)}      & $O(\log n)$ \\
Inspect the minimum     & \texttt{h[0]}            & $O(1)$ \\
Push then pop           & \texttt{heappushpop(h, x)} & $O(\log n)$, one pass \\
\bottomrule
\end{tabular}
\end{center}

A heap provides no ordering beyond its root. Iterating over the underlying
list does not produce sorted output.

\paragraph{Minimum only.} \texttt{heapq} implements a min-heap exclusively. A
max-heap is obtained by negating the keys on insertion and negating again on
extraction. For records, push the tuple \texttt{(-priority, item)}, since
tuples compare lexicographically and the first field dominates.

\begin{keybox}[Top-$k$ under a stream]
To retain the $k$ largest of $n$ elements, maintain a min-heap of size $k$.
Each new element is compared against the root, which is the smallest of the
current $k$. If the element is larger, the root is evicted and the element
inserted.

The cost is $O(n \log k)$ and the memory is $O(k)$. Sorting the full input
costs $O(n \log n)$ and requires all $n$ elements in memory at once, so the
heap is superior when $k \ll n$ and is the only option when the input arrives
as a stream of unknown length.
\end{keybox}

\begin{lstlisting}[style=python, caption={The k largest elements}]
import heapq

def top_k(stream, k):
    h = []                              # min-heap holding the k best so far
    for x in stream:
        if len(h) < k:
            heapq.heappush(h, x)
        elif x > h[0]:                  # h[0] is the weakest survivor
            heapq.heapreplace(h, x)     # pop the root, then push x
    return sorted(h, reverse=True)
\end{lstlisting}

When the whole input already fits in memory, the library provides this
directly as \texttt{heapq.nlargest(k, xs)} and \texttt{heapq.nsmallest(k, xs)},
both of which accept a \texttt{key} argument.

\paragraph{Two heaps for a running median.} Maintain a max-heap over the lower
half of the values and a min-heap over the upper half, rebalancing after each
insertion so their sizes differ by at most one. The median is then read from
the root of the larger heap, or is the mean of the two roots when the sizes
are equal. Insertion costs $O(\log n)$ and the median query is $O(1)$.
